% ============================================================
% Chapter 4: Proposed Method: Specification and Design
%   Section: The Fact-Checking Pipeline
%     Subsection: Verification with a Language Model
% Included from chapters/04-proposed-method/06-fact-checking-pipeline/fact-checking-pipeline.tex
% Source: docs/arquitectura-tecnica.md 2.4.5; verification/
% (per_source.py for the one-source-at-a-time rule).
% ============================================================

\subsection{Verification with a Language Model}
\label{subsec:llm-verification}

The model is given the claim, the article's headline and opening (so
that it reads the claim in its article's sense), and the ranked
evidence, numbered. It must return a structured answer:

\begin{itemize}
    \item a verdict from a closed set of five --- \texttt{TRUE},
    \texttt{PARTIALLY\_TRUE} (the central assertion holds, a detail does
    not), \texttt{MISLEADING} (technically true but deceptive, or
    sources of equal standing conflicting), \texttt{FALSE} and
    \texttt{UNVERIFIED};
    \item a confidence and a short explanation;
    \item the numbers of the evidence it used; and
    \item for each source, its stance --- \emph{supports},
    \emph{contradicts} or \emph{unrelated} --- with the verbatim passage
    it relied on.
\end{itemize}

It is told to answer \texttt{UNVERIFIED} when the evidence is empty or
does not bear on the claim, never to assert what the evidence does not
support, and to mark as \emph{unrelated} a source about a different
event even when it shares a place, a date or a figure with the claim.

Reducing the task to a closed classification with cited, quoted
evidence --- instead of a free-form verdict --- turns the question of
whether the model hallucinated into something that can be checked
mechanically in the next stage. Any quote that does not appear
verbatim in its source is discarded: an invented quotation is worse
than none. The model is served locally by default
(\texttt{llama3.2:3b}, a 3-billion-parameter model chosen for its
footprint, not by a benchmark), and any OpenAI-compatible provider can
take its place by configuration.

\paragraph{One source at a time.} Since 22~October the default asks
the model about each ranked source separately, in a short call: whether
that source supports the claim, contradicts it, is silent on its point
or is about something else, with the verbatim passage. The verdict then
follows a written rule rather than the model's judgement. Outlets are
counted, not pages --- five copies of one wire story are one source ---
and the side more outlets take wins; a tie is \texttt{UNVERIFIED}, and
so is a set of sources none of which takes a side. The confidence grows
with the share of outlets that agree, and the explanation is built from
the stances and their checked quotes, so that it cannot tell a different
story from them. This form never answers \texttt{PARTIALLY\_TRUE} or
\texttt{MISLEADING}. It was adopted on the results of
Section~\ref{sec:results-rq2}; the single call above remains a setting.

If the model cannot be reached after its retries, the claim is marked
\texttt{llmUnreachable} and reported apart from the verdicts, as a
failed search is. If there is no evidence at all, the model is not
asked: its answer would be discarded by the next stage anyway, and the
call is the slowest step of the pipeline.
